\subsection{DERIVATIVE OF A VECTOR FUNCTION}

DEFINITION 1. Let \(\mathbf{f}\) be a vector function defined on an interval \(\mathbf{I} \subset \mathbf{R}\) which does not reduce to a single point. We say that \(\mathbf{f}\) is differentiable at a point \(x_0 \in \mathbf{I}\) if

\(\lim_{x\to x_{0},x\in I,x\neq x_{0}}\frac{\mathbf{f}(x)-\mathbf{f}(x_{0})}{x-x_{0}}\) exists (in the vector space where f takes its values); the

value of this limit is called the first derivative (or simply the derivative) of f at the point \(x_{0}\) , and it is denoted by \(\mathbf{f}'(x_{0})\) or \(\mathbf{D}\mathbf{f}(x_{0})\) .

If f is differentiable at the point \(x_{0}\) , so is the restriction of f to any interval \(J \subset I\) which does not reduce to a single point and such that \(x_{0} \in J\) ; and the derivative of this restriction is equal to \(\mathbf{f}'(x_{0})\) . Conversely, let J be an interval contained in I and containing a neighbourhood of \(x_{0}\) relative to I; if the restriction of f to J admits a derivative at the point \(x_{0}\) , then so does f.

We summarise these properties by saying that the concept of derivative is a local concept.

Remarks. *1) In Kinematics, if the point \(\mathbf{f}(t)\) is the position of a moving point in the space \(\mathbf{R}^3\) at time \(t\) , then \(\frac{\mathbf{f}(t) - \mathbf{f}(t_0)}{t - t_0}\) is termed the average velocity between the instants \(t_0\) and \(t\) , and its limit \(\mathbf{f}'(t_0)\) is the instantaneous velocity (or simply velocity) at the time \(t_0\) (when this limit exists).*

\begin{enumerate}
\def\labelenumi{\arabic{enumi})}
\setcounter{enumi}{1}
\tightlist
\item
  If a function \(\mathbf{f}\) , defined on I, is differentiable at a point \(x_0 \in I\) , it is necessarily continuous relative to I at this point.
\end{enumerate}

DEFINITION 2. Let \(\mathbf{f}\) be a vector function defined on an interval \(\mathrm{I} \subset \mathbf{R}\) , and let \(x_0\) be a point of \(\mathrm{I}\) such that the interval \(\mathrm{I} \cap [x_0, +\infty[\) (resp. \(\mathrm{I} \cap ] - \infty, x_0]\) ) does not reduce to a single point. We say that \(\mathbf{f}\) is differentiable on the right (resp. on the left) at the point \(x_0\) if the restriction of \(\mathbf{f}\) to the interval \(\mathrm{I} \cap [x_0, +\infty[\) (resp. \(\mathrm{I} \cap ] - \infty, x_0]\) ) is differentiable at the point \(x_0\) ; the value of the derivative of this restriction at the point \(x_0\) is called the right (resp. left) derivative of \(\mathbf{f}\) at the point \(x_0\) and is denoted by \(\mathbf{f}_d'(x_0)\) (resp. \(\mathbf{f}_g'(x_0)\) ).

Let f be a vector function defined on I, and \(x_{0}\) an interior point of I such that f is continuous at this point; it follows from defs. 1 and 2 that for f to be differentiable at \(x_{0}\) it is necessary and sufficient that f admit both a right and a left derivative at this point, and that these derivatives be equal; and then

\[
\mathbf {f} ^ {\prime} (x _ {0}) = \mathbf {f} _ {d} ^ {\prime} (x _ {0}) = \mathbf {f} _ {g} ^ {\prime} (x _ {0}).
\]

Examples. 1) A constant function has zero derivative at every point.

\begin{enumerate}
\def\labelenumi{\arabic{enumi})}
\setcounter{enumi}{1}
\item
  An affine linear function \(x \mapsto ax + b\) has derivative equal to a at every point.
\item
  The real function \(1 / x\) (defined for \(x \neq 0\) ) is differentiable at each point \(x_0 \neq 0\) , for we have \(\left(\frac{1}{x} - \frac{1}{x_0}\right) / (x - x_0) = -\frac{1}{xx_0}\) , and, since \(1 / x\) is continuous at \(x_0\) , the limit of the preceding expression is \(-1 / x_0^2\) .
\item
  The scalar function \(|x|\) , defined on \(\mathbf{R}\) , has right derivative \(+1\) and left derivative \(-1\) at \(x = 0\) ; it is not differentiable at this point.
\end{enumerate}

*5) The real function equal to 0 for \(x = 0\) , and to \(x \sin 1/x\) for \(x \neq 0\) , is defined and continuous on \(\mathbf{R}\) , but has neither right nor left derivative at the point \(x \neq 0\) . One can give examples of functions which are continuous on an interval and fail to have a derivative at every point of the interval (I, p.~35, exerc. 2 and 3).

DEFINITION 3. We say that a vector function \(\mathbf{f}\) defined on an interval \(\mathrm{I} \subset \mathbf{R}\) is differentiable (resp. right differentiable, left differentiable) on \(\mathrm{I}\) if it is differentiable (resp. right differentiable, left differentiable) at each point of \(\mathrm{I}\) ; the function \(x \mapsto \mathbf{f}'(x)\) (resp. \(x \mapsto \mathbf{f}_d'(x)\) , \(x \mapsto \mathbf{f}_g'(x)\) ) defined on \(\mathrm{I}\) , is called the derived function, or (by abuse of language) the derivative (resp. right derivative, left derivative) of \(\mathbf{f}\) , and is denoted by \(\mathbf{f}'\) or \(\mathrm{Df}\) or \(df/dx\) (resp. \(\mathbf{f}_d', \mathbf{f}_g'\) ).

Remark. A function may be differentiable on an interval without its derivative being continuous at every point of the interval (cf.~I, p.~36, exerc. 5); *this is shown by the example of the function equal to 0 for x = 0 and to \(x^{2} \sin 1/x\) for \(x \neq 0\) ; it has a derivative everywhere, but this derivative is discontinuous at the point \(x = 0\)*

